\documentclass{beamer}

\usetheme[plain]{NTU}

\usepackage[absolute,overlay]{textpos}
\usepackage{tikz}
\usepackage[english]{babel}
\usepackage[latin1]{inputenc}
\usepackage{times}
\usepackage[T1]{fontenc}
\usepackage{tikz} % For drawing diagrams
\usetikzlibrary{shapes,arrows,positioning,automata,fit}
% Or whatever. Note that the encoding and the font should match. If T1
% does not look nice, try deleting the line with the fontenc.

\newcommand<>{\PutAt}[3][0pt]{%
    {\only#4{\begin{textblock*}{#1}#2%
      #3
    \end{textblock*}}}%
}

\newcommand{\ShowPutAtGrid}{
    \begin{textblock*}{128mm}(0cm,0cm)
    \tikz[color=red!20!white]\draw[very thin, step=5mm] (0mm,0mm) grid (130mm,100mm);
    \end{textblock*}
    \begin{textblock*}{128mm}(0cm,0cm)
    \begin{tikzpicture}[color=red]
      \draw[step=1cm] (0,0mm) grid (130mm,100mm);   
      \foreach \n in {0,...,12}
        \draw[xshift=.5mm,yshift=-1.5mm, inner sep=0pt, anchor=west] (\n,10) node {\scriptsize{\textbf{\n}}};
      \foreach \n in {1,...,9}
        \draw[xshift=.5mm,yshift=-1.5mm, inner sep=0pt, anchor=west] (0,10-\n) node {\scriptsize{\textbf{\n}}};
    \end{tikzpicture}
    \end{textblock*}
}

\newcommand<>{\NormalBox}[2][]{%
  \only#3{\tikz[#1, every node/.style={shape=rectangle,draw,fill=white, #1}]\node []{#2};}
}

\setbeamercolor{blue}{fg=blue!50!black}
\setbeamercolor{red}{fg=red!80!black}
\setbeamercolor{magenta}{fg=brown!50!red!60!black}

\newcommand{\blue}{blue!50!black} 
\newcommand{\red}{red!80!black} 
\newcommand{\green}{yellow!95!black} 
\newcommand{\magenta}{brown!50!red!60!black}

\newcommand*\oldmacro{}%
\let\oldmacro\insertshorttitle%
\renewcommand*\insertshorttitle{%
  \oldmacro\hfill%
  \insertframenumber\,/\,\inserttotalframenumber}


\title[]{Assured Machine Learning} % The short title appears at the bottom of every slide, the full title is only on the title page

%\subtitle
\author[]{Xiaozhe Gu, Arvind Easwaran}
\institute[NTU]{School of Computer Science and Engineering \\ Energy Research Institute (ERI@N) \\ \vspace{2mm} Nanyang Technological University, Singapore}



\date[June, 2018]{June, 2018}



% Delete this, if you do not want the table of contents to pop up at
% the beginning of each subsection:
\AtBeginSection[]
{
  \begin{frame}<beamer>{Outline}
    \tableofcontents[currentsection]
 \end{frame}
}


% If you wish to uncover everything in a step-wise fashion, uncomment
% the following command: 

% \beamerdefaultoverlayspecification{<+->}

% \setbeameroption{show notes on second screen=left}

\begin{document}

\begin{frame}
  \titlepage
\end{frame}

% \begin{frame}{Outline}
  % \tableofcontents
  % You might wish to add the option [pausesections]
% \end{frame}

% \section{System Model \& Motivation}




\begin{frame}[t]{Machine Learning  Applications in Safety-Critical Environments}
\begin{columns}[T] % align columns
\begin{column}[t]{.52\textwidth}
\vspace{-10mm}
\color{red}\rule{\linewidth}{2pt}
\vspace{-5mm}
\begin{itemize}
\visible<1,2,3>{
  \item Decision making in life-threatening conditions (Machine-Learning (ML) based medical decision support systems)}
  \visible<2,3>{
  \item  Robots (surgical robots, industrial robots, etc.)}
  \visible<3>{
    \item  Autonomous vehicles

\begin{figure}[h]
\centering
\caption{Autonomous Shuttle}
\includegraphics[scale=0.07]{fig/car.jpeg} 
\end{figure}
}
\end{itemize}

\end{column}%
\hfill%


\begin{column}[t]{.44\textwidth}
\vspace{-10mm}
\color{blue}\rule{\linewidth}{2pt}
\vspace{-5mm}

\visible<1,2,3>{
\begin{figure}[h]
\centering
\caption{ML Based Brain Disease Diagnosis$^1$}
\includegraphics[scale=0.2]{fig/medical.png} 
\end{figure}
}

\vspace{-10mm}
\visible<2,3>{
 \begin{figure}[h]
\centering
\caption{Surgical Robots$^2$}
\includegraphics[scale=0.08]{fig/surgical.jpg} 
\end{figure}
}




\end{column}%
\end{columns}

  
\end{frame}



\begin{frame}[t]{Challenges to Safety Assurance}

\setbeamercovered{transparent}
\begin{itemize}
\visible<1,2,3,4>{
  \item  \textcolor{blue}{Non-transparency}: It is difficult to assess the reliability if  the reasoning behind  these models cannot be understood
\vspace{2mm}
  }
  \visible<2,3,4>{
  \item  \textcolor{blue}{Error Rate}:The estimate of error rate of a ML model with respect to the test data is not reliable
\vspace{2mm}
  }
  \visible<3,4>{
  \item     \textcolor{blue}{Instability}: A small change in the training process may produce a different result, and hence it is difficult to debug models or reuse parts of previous safety assessments.
\vspace{2mm}
  }
  \visible<4>{
  \item \textcolor{blue}{Difficulty in verification}: Formal verification of ML components is a difficult, and somewhat ill-posed, problem due to the complexity of the underlying ML algorithms and large feature spaces
}
\end{itemize}

\end{frame}



\begin{frame}[t]{Potential Strategies for Safety Assurance}
\begin{columns}[T] % align columns
\begin{column}[t]{.56\textwidth}
\vspace{-10mm}
\color{red}\rule{\linewidth}{2pt}
\vspace{-5mm}
\begin{itemize}
\visible<1,2,3>{
  \item  \textcolor{blue}{Interpretability \& Transparency}: Improve the  interpretability \& transparency of the ML component}
  \visible<2,3>{
  \item   \textcolor{blue}{Fail-Safe}: The model reports that it cannot reliably give a prediction and does not attempt to do so, thereby failing safely 
  }
  \visible<3>{
    \item \textcolor{blue}{Abstract}: Abstract the ML component and input feature space, and identify scenarios that could cause violation of safety specifications
}
\end{itemize}

\end{column}%
\hfill%


\begin{column}[t]{.4\textwidth}
\vspace{-10mm}
\color{blue}\rule{\linewidth}{2pt}
\vspace{-10mm}
\visible<1,2,3>{
\begin{figure}[h]
\centering
\caption{Explanations improve trust in prediction~\cite{lime}}
\includegraphics[scale=0.25]{fig/lime.png} 
\end{figure}
}


\vspace{-2mm}
\visible<2,3>{
\color{black}
Technique used in ML when predictions cannot be given confidently is the {\color{red} reject option}~\cite{reject}

\vspace{-3mm}
\begin{small}
  \begin{align*}
  \hat y(x)=\begin{cases}-1~\mbox{ if  }\phi(x)\leq t\\\mbox{reject, if }\phi( x)\in(-t,t)\\1~\mbox{ if  }\phi(x)\geq t\end{cases}
\end{align*}
\end{small}


}




\end{column}%
\end{columns}

  
\end{frame}





\renewcommand{\refname}{~}
\begin{thebibliography}{1}

\bibitem{}http://mlcenter.postech.ac.kr/healthcare
\vspace{2mm}


\bibitem{}https://www.wired.com/2015/03/google-robot-surgery/
\vspace{2mm}

\bibitem{lime} Ribeiro, Marco Tulio, Sameer Singh, and Carlos Guestrin. "Why should i trust you?: Explaining the predictions of any classifier." Proceedings of the 22nd ACM SIGKDD International Conference on Knowledge Discovery and Data Mining. ACM, 2016.
\vspace{2mm}

\bibitem{reject} Bartlett P L, Wegkamp M H. Classification with a reject option using a hinge loss[J]. Journal of Machine Learning Research, 2008, 9(Aug): 1823-1840.

\end{thebibliography}

\end{document}